\subsection{Puncture recursion and exact genus-two profinite ranks}

\begin{result}{D.37}{Proved}{The surviving-complex comparison and all genus-two ranks}
Put
\[
 \mathcal Y_n=\mathcal C_{2,n}^{\rm lev}
              =\check{\mathcal C}_{2,n}.
\]
Here the curve complex is nonempty, and the equality is the actual
full-to-congruence comparison of D.19.  For $n\geq1$, let $Z_n$
be the full subcomplex of curves that remain essential on forgetting
the last puncture.  Give its completion $Z_n^{\rm lev}$ the levels
induced by $\Gamma_{2,n}$.  Actual forgetting induces, for every
finite constant abelian coefficient group $A$, isomorphisms
\[
 r^*:H^j_{\rm cts}(\mathcal Y_{n-1};A)
       \ \xrightarrow{\ \sim\ }\
       H^j_{\rm cts}(Z_n^{\rm lev};A)
       \qquad(j\geq0).                                      \tag{D.69}
\]
For $n=1$ there are no deleted vertices, so
$Z_1^{\rm lev}=\mathcal Y_1$.

For $n\geq2$, let $B_n$ be the profinite set of vertices deleted
by this forgetting: they bound a disk containing the last puncture
and one other puncture.  For every $j\geq2$ there is a natural
isomorphism
\[
 H^j_{\rm cts}(\mathcal Y_n;A)
   \ \cong\
 \frac{C_{\rm cont}\bigl(B_n,
             H^{j-1}_{\rm cts}(\mathcal Y_{n-1};A)\bigr)}
      {H^{j-1}_{\rm cts}(\mathcal Y_{n-1};A)}.                 \tag{D.70}
\]
The denominator is the subgroup of constant functions.
The cohomology group in the target of $C_{\rm cont}$ is given the
discrete topology; these are locally constant functions with finite
image, equivalently the union over finite clopen partitions of
$B_n$.

Set $h_0=2$ and $h_n=n+1$ for $n\geq1$.  Positive-degree
continuous cohomology of $\mathcal Y_n$ with any finite constant
coefficients is zero outside degree $h_n$, and
\[
       H^{h_n}_{\rm cts}(\mathcal Y_n;\mathbb F_2)\ne0 .
                                                               \tag{D.71}
\]
Consequently the full and procongruence completions of every
genus-two rank have exactly the discrete connectivities of D.36:
for $k<d=3+n$, they are $-1,0,1$ at $k=0,1,2$ respectively,
and $k-2$ for $3\leq k<d$.  For $k\geq d$ they are contractible.
Connectivity and exactness have the conventions preceding D.35.
\end{result}
\begin{proof}
We give the actual maps in both comparisons; in particular no
commutation of arbitrary inverse limits with pushouts or equalizers
is assumed.

Write $G=\Gamma_{2,n}$, $Q=\Gamma_{2,n-1}$ and
$P=\pi_1(S_{2,n-1})$.  The actual forgetting sequence is
$1\to P\to G\to Q\to1$, with its full completed exact sequence
and full topology on $P$ as in D.31.  Kent--Leininger--Schleimer's
surviving-complex theorem identifies the actual map
$r:Z_n\to\mathcal C_{2,n-1}$ as a homotopy equivalence in the
polyhedral topology
\cite[Theorem 7.1, Corollary 7.2, and Proposition 7.6]
{KentLeiningerSchleimer2009}.
The opened proof gives tree fibres and the contraction along
their unique arcs; Proposition 7.6 proves continuity in precisely
that topology.  This is the CW topology used for the Borel
constructions.  The map is $G$-equivariant, with the target action
through $Q$.

For a finite-index normal subgroup $H$ of $G$, put
$J=q(H)$ and $N=H\cap P$.  Equivariant forgetting gives
the first map in
\[
 EH\times_H Z_n
       \longrightarrow EH\times_H\mathcal C_{2,n-1}
       \longrightarrow EJ\times_J\mathcal C_{2,n-1}.           \tag{D.72}
\]
The first is a homotopy equivalence: it is a map of Borel
fibrations over $BH$ and is the stated equivalence on their
fibres.  The second has homotopy fibre $BN$.  More explicitly,
it is the pullback, in the homotopy sense, of $BH\to BJ$ along
$EJ\times_J\mathcal C_{2,n-1}\to BJ$; the extension
$1\to N\to H\to J\to1$ identifies that fibre.

We verify the passage to levels in all degrees.  The orbit-to-Borel
comparison of (D.65) applies to $Z_n$ and to
$\mathcal C_{2,n-1}$ without a connectivity restriction.  Their
cell stabilizers are the actual curve or flag stabilizers of
genus-two groups, hence are good and have full ambient topology
by the proof of D.35.  There are finitely many cell orbits in each
degree and a fixed finite dimension.  Its same normal-refinement
argument therefore gives
\[
\begin{split}
 H^*_{\rm cts}(Z_n^{\rm lev};A)
   &\cong\varinjlim_H H^*(EH\times_H Z_n;A),\\
 H^*_{\rm cts}(\mathcal Y_{n-1};A)
   &\cong\varinjlim_J
             H^*(EJ\times_J\mathcal C_{2,n-1};A).
\end{split}                                                    \tag{D.73}
\]
All arrows here are induced by forgetting the contractible
coordinate, as in (D.64).

The second map of (D.72) has the cohomology spectral sequence
with local coefficients $H^q(N;A)$.  These groups are finite:
$N$ is a finite-index free group when $n\geq2$, and a
finite-index closed surface group when $n=1$.  They vanish above
degree one or two, respectively.  Both kinds of group are good.
For a fixed $H$, kill every element of every positive-degree
$H^q(N;A)$ on one finite-index subgroup of $N$.  This is possible
by goodness and finiteness, exactly as in D.35.  Full topology
of $P$ in $G$ then gives $H'\triangleleft_f G$, $H'\subset H$,
such that $H'\cap P$ lies in that subgroup.  Thus the map on
the entire positive-degree coefficient group is zero.  It
follows that the induced map on every spectral-sequence term
in a positive fibre degree is zero, even when the corresponding
base cohomology group is infinite.  No finiteness of that base
cohomology is required.  The row $q=0$ is the ordinary constant
coefficient cohomology of the target of (D.72).

The subgroups $q(H)$ are cofinal among the finite-index normal
subgroups of $Q$: the inverse image of every such subgroup is
an allowed $H$.  Exactness of filtered colimits and the bounded
fibre degree now show that the second map in (D.72) induces an
isomorphism after these colimits.  Combining with the first map
and (D.73) proves (D.69).  Every map in this argument is induced
by the original forgetting, so this is precisely $r^*$.
For $n=1$, a deleted essential curve would have to bound a disk
with the sole puncture, which is already inessential.  Hence
there is no deletion in that case.

We next identify the completed link bundle.  Distinct deleted
curves cannot be disjoint.  Their twice-punctured disks both
contain the forgotten puncture; disjoint boundaries force the
disks to be nested.  The enclosed other puncture is then the
same, and the annulus between their boundaries has no punctures,
so the two curves are isotopic.  This proves the assertion also
when the other labels are specified.  It follows that every
simplex either belongs to $Z_n$ or contains exactly one deleted
vertex.

Let $L_n$ consist of pairs $(b,\sigma)$ where $b$ is a deleted
vertex and $\sigma$ is a nonempty simplex in its link.  Forgetting
identifies this actual link with $\mathcal C_{2,n-1}$: cut off
the twice-marked disk and identify its boundary collar with the
retained puncture.  We prove that the same map on levels is the
homeomorphism
\[
 L_n^{\rm lev}\longrightarrow B_n\times\mathcal Y_{n-1},
       \qquad (b,\sigma)\longmapsto(b,r(\sigma)).
                                                               \tag{D.74}
\]
Fix a deleted curve $b$ of one of its finitely many labelled
types.  Its stabilizer $C_b$ has the actual cutting sequence
$1\to\mathbb Z_b\to C_b\to Q\to1$.  The boundary twist is its
kernel, and every mapping class of the core lifts.  D.19 and
the valid cutting theorem, including its link identification
in Boggi's Remark 4.7, give the exact completed sequence with
full topology on $C_b$.  If $\sigma$ is a simplex of the core
with stabilizer $S_\sigma$ in $Q$, then $C_{b,\sigma}$ is exactly
the inverse image of $S_\sigma$.  Its completion is the inverse
image of $\widehat S_\sigma$, by the same exact cutting
comparison.  Consequently the map on each orbit is
\[
 \widehat G/\widehat C_{b,\sigma}
 \longrightarrow
   (\widehat G/\widehat C_b)\times
   (\widehat Q/\widehat S_\sigma),\qquad
 g\longmapsto(g\widehat C_b,\widehat q(g)\widehat S_\sigma).
                                                               \tag{D.75}
\]
It is surjective because $\widehat C_b\to\widehat Q$ is
surjective: after choosing the first coset, multiplication by
an element of $\widehat C_b$ adjusts the second to any desired
coset.  It is injective because
$\widehat C_{b,\sigma}$ is precisely the displayed inverse
image inside $\widehat C_b$.  This continuous bijection is a
homeomorphism by compactness.  The finitely many simplex
types, and the face maps obtained by forgetting members of
$\sigma$, assemble these maps into (D.74).  This proves the
bundle assertion using the actual forgetful map and does not
choose a splitting of the central cutting extension.

The set $B_n$ is a clopen union of vertex-type orbits; the
surviving simplices likewise form clopen pieces in each degree.
The statement that there is at most one deleted vertex persists
at every level.  A simplex in the other piece has a unique
deleted vertex and a unique opposite surviving face.  Applying
these two assignments and their inverses in each degree proves
the following actual cellular pushout after completion:
\[
 \mathcal Y_n =
 Z_n^{\rm lev}\ \underset{L_n^{\rm lev}}{\cup}\
             \operatorname{Cone}_{B_n}(L_n^{\rm lev}).
                                                               \tag{D.76}
\]
The cone is taken separately over each $b\in B_n$; its cone
vertices give the isolated deleted vertices as well.
One may barycentrically subdivide these cells before taking
levels, so that stabilizers have no inversions.  The explicit
unique-vertex and opposite-face assignments prove (D.76) for
that model too.  Thus (D.76) is not inferred from a general
inverse-limit preservation claim.

The fibrewise cone retracts to the profinite set $B_n$.
The continuous-cochain Mayer--Vietoris sequence applies to
(D.76).  Concretely its cochain short exact sequence is the
usual difference of the two restrictions to $L_n^{\rm lev}$.
Surjectivity follows from the closed base inclusion in the
fibrewise cone: a continuous finite-valued function on a
closed subspace of a profinite set extends, by a finite clopen
refinement.  This argument does not require the other map
$L_n^{\rm lev}\to Z_n^{\rm lev}$ to be injective.

For any profinite $B$ and simplicial profinite $Y$, the
cochain complex of $B\times Y$ with finite constant
coefficients is $C_{\rm cont}(B,C^*_{\rm cts}(Y;A))$,
where the cochain groups inside are discrete.  Indeed a
finite-valued continuous function has a finite rectangular
clopen refinement.  The functor of locally constant
finite-image functions on $B$ is exact: for a surjection,
choose preimages of the finitely many values of a given
function.  Therefore (D.74) gives
\[
 H^j_{\rm cts}(L_n^{\rm lev};A)
  =C_{\rm cont}\bigl(B_n,
                     H^j_{\rm cts}(\mathcal Y_{n-1};A)\bigr).
                                                               \tag{D.77}
\]
Under (D.69), the map from $H^j(Z_n^{\rm lev};A)$ to (D.77)
is the constant-function inclusion: restriction to every
link is the same actual forgetting map in (D.74).  In
positive degrees the cohomology of the cone term is zero,
and the constant-function inclusion is injective since
$B_n$ is nonempty.  The Mayer--Vietoris sequence in degrees
$j-1,j$ proves exactly (D.70).

The base $\mathcal Y_0$ has dimension two, has trivial
fundamental group by D.11, and has nonzero degree-two
$\mathbb F_2$ cohomology by the separated finite-cycle
argument of D.35.  Hence its positive finite-coefficient
cohomology is concentrated in degree two.  Equation (D.69)
gives the identical assertion for $\mathcal Y_1$.
For $n\geq2$, the set $B_n$ has at least two distinct points.
For example a twice-punctured disk curve and its image by
a point push around a handle are distinct curves of the
same deleted type.  The actual curve embedding and
congruence separability of D.35 separate them at a finite
level.  They therefore admit a continuous function to
$\{0,1\}$ taking both values.  Multiplying such a function
by any nonzero cohomology class of the preceding stage
produces a nonconstant element of the numerator of
(D.70).  Thus the recursion shifts the sole possible
positive cohomology degree by one and preserves
nonvanishing there with $\mathbb F_2$ coefficients.
The vanishing of $H^1$ follows also from D.11.
This proves (D.71) and the concentration assertion.

For completeness we give the continuous maps that detect
the remaining rank-selected groups.  The sections in
D.7 and D.36 fix a rational-tail multicurve and adjoin
its curves to a multicurve on the complementary core.
The valid procongruence link theorem, Boggi Remark 4.7
with the actual cutting map of Theorem 4.5, identifies
that core completion with the corresponding closed
subcomplex of the profinite link.  Union with the fixed
tail curves is a continuous map of face posets.  Ordinary
forgetting is continuous at full levels, since the
inverse image of a finite-index subgroup is finite
index.  The edge-count calculation of D.36 ensures it
has precisely the indicated rank-selected target.
The two composites are the identity on the dense
actual face posets, hence on their completions.  D.19
identifies the full and congruence versions of these
maps.  This proves continuous retractions, without
requiring a section of a central framing extension.

For $3\leq k\leq d-2$, their composite is explicitly
\[
 P_k(2,n)^{\rm lev}\longrightarrow
 P_k(2,k-1)^{\rm lev}\longrightarrow
 \mathcal C_{2,k-2}^{\rm lev},                                \tag{D.78}
\]
with the nested-tail section first and the one-tail
section second in the reverse direction.  Equation
(D.71) gives nonzero cohomology in degree $k-1$ on
the last space, which pulls back nontrivially through
this retraction.  D.35 gives the lower connectivity
$k-2$.  At $k=d-1$ use (D.71) directly.  At $k=0,1,2$,
the finite top-dimensional cycles constructed for
every genus-two type in D.36 remain separated at one
finite level and are detected exactly as in D.35.
The lower bounds there are $-1,0,1$.
To infer nonzero homotopy from the nonzero cohomology,
use the proved lower connectivity, the coefficient
sequence into $\mathbb Q/\mathbb Z$, Pontryagin duality
and profinite Hurewicz exactly as in D.35.  This proves
the full list of exact connectivities.
\end{proof}

The higher-genus groups are not evaluated by this recursion.
In particular D.32 still identifies
$\Pi_1(\check P_2(3,0))$ with the unevaluated congruence kernel
$K_{3,0}$.  D.34 gives the exact kernel quotients at all
other ranks, and (D.63) gives the stated full-level
cohomology groups as ambient restriction colimits;
their general higher-genus vanishing is not proved here.
